\section{Jenis Pengujian: Fungsional, Kompatibilitas, dan Performa}

Pengujian aplikasi web dibagi menjadi beberapa jenis berdasarkan aspek yang diuji. \textbf{Pengujian fungsional} memeriksa apakah fitur bekerja sesuai spesifikasi --- misalnya form registrasi menyimpan data, login berhasil dengan kredensial benar, dan menolak yang salah. \textbf{Pengujian kompatibilitas} memastikan aplikasi tampil dan berfungsi dengan baik di berbagai browser (Chrome, Firefox, Safari, Edge), sistem operasi, dan ukuran layar (responsive design). \textbf{Pengujian performa} mengukur kecepatan, kapasitas, dan responsivitas aplikasi di bawah berbagai beban \cite{mdnToolsTesting}.

\begin{table}[ht]
\centering
\begin{tabular}{|P{2.5cm}|P{4.5cm}|P{3cm}|P{2.5cm}|}
\hline
\textbf{Jenis} & \textbf{Fokus} & \textbf{Alat} & \textbf{Kapan} \\
\hline
Fungsional & Fitur sesuai spesifikasi & Selenium, Cypress & Setiap build \\
\hline
Kompatibilitas & Browser, OS, device & BrowserStack & Sebelum rilis \\
\hline
Performa & Kecepatan, responsivitas & Lighthouse, k6 & Sebelum deploy \\
\hline
Keamanan & Celah dan kerentanan & OWASP ZAP & Berkala \\
\hline
Usabilitas & Kemudahan penggunaan & User testing & Desain awal \\
\hline
\end{tabular}
\caption{Jenis-jenis pengujian aplikasi web}
\end{table}

